% Chapter 8: Tornado Cash on Cardano: A Case Study.
% Academic analysis only. No deployable implementation is provided.

%--------------------------------------------------------------------------
\section{Scope and Purpose} \label{sec:ch8-scope}

\begin{highlight}
\textbf{This chapter is an academic case study.} Its purpose is to take the most
demanding well-known ZK application and work out, concretely, whether Cardano
can host it today. The answer turns out to be \emph{no}, and the reasons why
teach more about the platform's real limits than any working example could.

No deployable implementation is provided. Nothing in this chapter is deployed to
any network, and no complete circuit or validator set is given. The code
fragments are design sketches for analysis.
\end{highlight}

There is a second reason for the disclaimer. Tornado Cash was sanctioned by the
US Office of Foreign Assets Control in August 2022, and a developer was
subsequently prosecuted over its operation. Whatever one's view of those
actions, which were and remain contested, deploying a mixer carries
legal exposure that has nothing to do with whether the cryptography works. This
chapter analyses a design. It does not recommend building one.

What we are actually after is the engineering question, which is genuinely
interesting and has a genuinely surprising answer:

\begin{outline}
Cardano can verify a zero-knowledge proof for about 0.15 ADA and 20\% of a
transaction's CPU budget, at a cost that does not depend on the circuit's
complexity. Why, then, can it not host an application whose on-chain component
is one proof verification?
\end{outline}

Because verifying a proof and \emph{maintaining the state the proof is about}
are completely different problems, and Cardano is excellent at the first and
poor at the second. This chapter measures exactly how poor.

%--------------------------------------------------------------------------
\section{What Tornado Cash Actually Does} \label{sec:ch8-what}

Strip away the politics and Tornado Cash is a clean cryptographic construction:
a \textbf{shielded pool} built from three ingredients.

\subsection*{The Deposit}

A depositor generates two random 248-bit values, a \texttt{nullifier} and a
\texttt{secret}, and computes a \textbf{commitment}:

\[
  C = H(\texttt{nullifier}, \texttt{secret})
\]

They send a fixed amount (1 ETH, or 10, or 100) to the pool contract along
with $C$. The contract appends $C$ as a leaf to an append-only
\textbf{incremental Merkle tree} of depth 20 and stores the new root. The pair
$(\texttt{nullifier}, \texttt{secret})$ is the depositor's \emph{note}; losing it
loses the funds.

Crucially, the contract does not record who deposited which commitment beyond
what the transaction itself reveals. The commitment is just a hash.

\subsection*{The Withdrawal}

Later, and ideally much later, the depositor proves in zero knowledge that:

\begin{enumerate}
  \item they know a $(\texttt{nullifier}, \texttt{secret})$ pair whose
        commitment $C$ is a leaf of the tree with root $R$, \emph{without
        revealing which leaf}; and
  \item the public value $\texttt{nullifierHash} = H(\texttt{nullifier})$
        is correctly derived from that same nullifier.
\end{enumerate}

The contract checks the proof, checks that $R$ is a recent known root, checks
that \texttt{nullifierHash} has not been seen before, records it, and pays the
denomination to a recipient address supplied as a public input.

\subsection*{Why the Three Pieces Are Each Necessary}

\begin{center}
\begin{tabular}{lp{7.6cm}}
  \hline
  \textbf{Piece} & \textbf{What breaks without it} \\
  \hline
  Merkle tree        & No way to prove membership in the deposit set without
                       naming your deposit; the link is exposed. \\
  \addlinespace
  Nullifier          & The same note could be withdrawn repeatedly; the pool is
                       drained. \\
  \addlinespace
  Fixed denomination & Amounts become fingerprints. A deposit and a
                       withdrawal of 3.7194 ADA are trivially linked. \\
  \hline
\end{tabular}
\end{center}

The anonymity set is the number of deposits in the tree that could plausibly be
yours. A pool with 3 deposits provides essentially no privacy; Tornado's largest
pools held tens of thousands. \textbf{Anonymity is a function of pool
throughput}, and this will matter enormously later.

\subsection*{The Relayer}

There is a bootstrapping problem: to withdraw, you need a transaction, and a
transaction needs a fee, and a fee needs funds, and if you pay the fee from an
address linked to your deposit, you have undone the entire construction. Tornado
solves this with \textbf{relayers}: third parties who submit the withdrawal
transaction and pay the fee, taking a cut from the withdrawn amount. The relayer
learns the proof and the recipient but never the note.

%--------------------------------------------------------------------------
\section{Sketching It in eUTxO} \label{sec:ch8-sketch}

Let us design it for Cardano properly, giving the eUTxO model every advantage,
before looking for problems.

\subsection*{The Pool State}

Ethereum's contract has mutable storage. Cardano has UTxOs. The natural
translation is a single \textbf{pool UTxO} carrying the tree state in its datum,
consumed and recreated by every deposit:

\begin{verbatim}
pub type PoolDatum {
  root:       ByteArray,        // current Merkle root
  next_index: Int,              // number of leaves so far
  filled:     List<ByteArray>,  // rightmost node per level
  denom:      Int,              // fixed denomination, lovelace
}
\end{verbatim}

The \texttt{filled} list is the standard incremental-Merkle-tree trick: to
append a leaf and recompute the root you do not need the whole tree, only the
rightmost filled node at each of the 20 levels, plus a table of precomputed
"zero" subtree hashes. That is 20 hashes per insertion instead of a million.

\subsection*{The Nullifier Set: Where Cardano Wins}

Ethereum stores spent nullifiers in a \texttt{mapping(bytes32 => bool)}. On
Cardano, tracking a growing set is normally painful. But there is an unusually
elegant option here, and it is worth pausing on because it is genuinely better
than the Ethereum version.

\textbf{Mint a token whose asset name is the nullifier hash.}

The Cardano ledger enforces, as a consensus rule, that a given
(policy id, asset name) pair with quantity 1 cannot be minted twice under a
policy that only ever mints one. So the minting policy becomes the
double-spend check, and it is enforced by the ledger rather than by our script:

\begin{verbatim}
validator nullifier_policy {
  mint(redeemer: WithdrawRedeemer, policy_id: PolicyId,
       tx: Transaction) {
    expect [Pair(asset_name, 1)] =
      tx.mint |> assets.tokens(policy_id) |> dict.to_pairs()

    // The minted asset name IS the nullifier hash, so uniqueness
    // is a ledger invariant, not something we have to check.
    and {
      asset_name == redeemer.nullifier_hash,
      verify_withdraw_proof(redeemer),
    }
  }
}
\end{verbatim}

\begin{highlight}
This is a real advantage of the eUTxO model. On Ethereum, double-withdrawal
prevention costs an \texttt{SSTORE} and depends on the contract being correct.
On Cardano it costs a token mint and depends on the \emph{ledger} being correct.
No script bug can permit a double withdrawal.
\end{highlight}

So far Cardano is doing well. The trouble is elsewhere.

%--------------------------------------------------------------------------
\section{The Withdrawal Circuit} \label{sec:ch8-circuit}

The circuit is a direct adaptation of Tornado's, with Poseidon in place of the
original Pedersen/MiMC pair; Poseidon is what circomlib gives us over
BLS12-381, and it is what Chapter~\ref{ch:password} already used.

\begin{verbatim}
pragma circom 2.1.6;

include "../node_modules/circomlib/circuits/poseidon.circom";
include "../node_modules/circomlib/circuits/bitify.circom";

// One level of a Merkle path. `sel` is the path bit: 0 means our
// node is on the left, 1 means on the right.
template MerkleLevel() {
    signal input  cur;
    signal input  sibling;
    signal input  sel;
    signal output out;

    sel * (sel - 1) === 0;   // force sel to be a bit

    signal left;
    signal right;
    left  <== cur + sel * (sibling - cur);
    right <== sibling + sel * (cur - sibling);

    component h = Poseidon(2);
    h.inputs[0] <== left;
    h.inputs[1] <== right;
    out <== h.out;
}

template Withdraw(levels) {
    // ---- public inputs (checked by the Aiken validator) ----
    signal input root;
    signal input nullifierHash;
    signal input recipient;  // payment credential hash
    signal input fee;        // relayer fee in lovelace

    // ---- private witness (never leaves the prover) ----
    signal input nullifier;
    signal input secret;
    signal input pathElements[levels];
    signal input pathIndices[levels];

    // 1. commitment = Poseidon(nullifier, secret)
    component commit = Poseidon(2);
    commit.inputs[0] <== nullifier;
    commit.inputs[1] <== secret;

    // 2. nullifierHash = Poseidon(nullifier)
    component nh = Poseidon(1);
    nh.inputs[0] <== nullifier;
    nh.out === nullifierHash;

    // 3. the commitment is a leaf under `root`
    component path[levels];
    signal cur[levels + 1];
    cur[0] <== commit.out;

    for (var i = 0; i < levels; i++) {
        path[i] = MerkleLevel();
        path[i].cur     <== cur[i];
        path[i].sibling <== pathElements[i];
        path[i].sel     <== pathIndices[i];
        cur[i + 1]      <== path[i].out;
    }

    cur[levels] === root;

    // 4. Bind recipient and fee into the constraint system so a
    //    relayer cannot redirect the payment. Squaring is the
    //    standard Circom idiom for "this signal must be used".
    signal recipientSq;
    signal feeSq;
    recipientSq <== recipient * recipient;
    feeSq       <== fee * fee;
}

component main {public [root, nullifierHash, recipient, fee]}
    = Withdraw(20);
\end{verbatim}

Two Cardano-specific notes. First, \texttt{recipient} is a payment credential
hash (28 bytes, 224 bits), which fits comfortably in a 255-bit BLS12-381
field element, so the address can be a single public input. Second, binding
\texttt{recipient} and \texttt{fee} into the circuit is the replay protection
from Section~\ref{sec:ch7-ops}: a stolen proof can only ever pay the original
recipient.

This circuit is correct. Compiled for BLS12-381 it produces a proof that
Cardano's builtins can verify in 2.0 billion CPU steps, exactly as in
Chapter~\ref{ch:password}. \textbf{The withdrawal side of Tornado Cash is not the
problem.}

%--------------------------------------------------------------------------
\section{Wall One: There Is No On-Chain Poseidon} \label{sec:ch8-wall1}

The problem is the \emph{deposit}. To append a leaf, the pool validator must
recompute the Merkle root, and that means computing Poseidon \emph{on-chain},
in Aiken, 20 times.

Cardano has no Poseidon builtin. The available hash builtins are
\texttt{sha2\_256}, \texttt{sha3\_256}, \texttt{blake2b\_224},
\texttt{blake2b\_256}, \texttt{keccak\_256} and \texttt{ripemd\_160}; all
excellent, none of them ZK-friendly. So Poseidon would have to be written in
Aiken, from field arithmetic.

\subsection*{What Field Arithmetic Costs, Measured}

Aiken's \texttt{Scalar} type is not a native field element. Looking at the
stdlib source:

\begin{verbatim}
pub fn mul(left: Scalar, right: Scalar) -> Scalar {
  Scalar(left.integer * right.integer % field_prime)
}
\end{verbatim}

That is arbitrary-precision integer multiplication followed by a 255-bit
modulo, interpreted by the Plutus CEK machine. There is no BLS12-381 scalar
arithmetic builtin. So let us measure it rather than guess.

The benchmark file in the repository,
\texttt{lib/zk\_password/bench\_onchain.ak}, runs tail-recursive loops of
1000 iterations with matching \texttt{nop} baselines, so subtracting the
baseline and dividing by 1000 gives the marginal cost of one operation:

\begin{verbatim}
aiken check -m bench
\end{verbatim}

\begin{center}
\begin{tabular}{lrr}
  \hline
  \textbf{Operation} & \textbf{CPU steps} & \textbf{Mem units} \\
  \hline
  \texttt{acc + 3} (small ints)          & 197{,}240 & 602 \\
  \texttt{acc \% p}                      & 250{,}615 & 604 \\
  \texttt{scalar.add}                    & 417{,}083 & 1{,}009 \\
  \texttt{scalar.mul}                    & 421{,}644 & 1{,}011 \\
  raw \texttt{(a + b) \% p}              & 417{,}083 & 1{,}009 \\
  raw \texttt{a * b \% p}                & 430{,}929 & 1{,}012 \\
  \hline
  \texttt{blake2b\_256} (64-byte input)  & 398{,}720 & 813 \\
  \texttt{sha2\_256} (64-byte input)     & 581{,}923 & 813 \\
  \hline
\end{tabular}
\end{center}

Three things in this table matter more than the headline numbers.

\begin{enumerate}
  \item \textbf{The \texttt{Scalar} wrapper is not the problem.} Raw
        \texttt{a * b \% p} on plain integers costs the same as
        \texttt{scalar.mul}. Bypassing the stdlib buys nothing.
  \item \textbf{Operand width is not the problem either.} A bare small-integer
        \texttt{acc + 3} already costs 197{,}240 CPU steps. The 255-bit modular
        multiply costs only about twice that.
  \item \textbf{Therefore the cost is the interpreter, not the arithmetic.}
        Roughly 200{,}000 CPU steps and 600 memory units are the price of
        \emph{any} elementary operation in a Plutus script. The numbers are
        dominated by CEK machine overhead per reduction step.
\end{enumerate}

\begin{highlight}
This is the most consequential fact in the chapter. Because the cost floor is
per-\emph{operation} rather than per-\emph{bit}, no optimisation helps.
Montgomery form, lazy reduction, batched inversion, a hand-tuned UPLC
implementation; all of them reduce the \emph{width} of the arithmetic, and
width is not what you are paying for. You are paying about $2 \times 10^5$ CPU
steps per operation, and Poseidon needs about $1.4 \times 10^3$ of them.
\end{highlight}

\subsection*{Poseidon, Costed}

Poseidon over BLS12-381 with $t = 3$ (the 2-to-1 compression a binary Merkle
tree needs) uses circomlib's standard parameters: 8 full rounds and 57 partial
rounds, with an $x^5$ S-box.

\begin{center}
\begin{tabular}{lrrr}
  \hline
  \textbf{Per round} & \textbf{Mults} & \textbf{Adds} & \textbf{Rounds} \\
  \hline
  Full: 3 S-boxes ($3\times$), MDS ($9\times$) & 18 & 9 & 8 \\
  Partial: 1 S-box ($3\times$), MDS ($9\times$) & 12 & 9 & 57 \\
  \hline
  \textbf{Total} & \textbf{828} & \textbf{585} & \\
  \hline
\end{tabular}
\end{center}

At the measured rates:

\[
  828 \times 430{,}929 \;+\; 585 \times 417{,}083
  \;=\; 600{,}802{,}860 \text{ CPU steps}
\]
\[
  828 \times 1{,}012 \;+\; 585 \times 1{,}009
  \;=\; 1{,}428{,}386 \text{ memory units}
\]

So \textbf{one} on-chain Poseidon compression costs 6.0\% of the transaction CPU
budget and 8.7\% of the memory budget. About 0.126 ADA per hash.

\subsection*{The Verdict on the Deposit}

A depth-20 insertion needs 20 of them:

\begin{center}
\begin{tabular}{lrrr}
  \hline
  \textbf{Depth-20 Merkle insertion} & \textbf{CPU} & \textbf{Memory} \\
  \hline
  Poseidon in Aiken (estimated) & 12{,}016{,}057{,}195 & 28{,}567{,}718 \\
  Transaction budget            & 10{,}000{,}000{,}000 & 16{,}500{,}000 \\
  \textbf{Utilisation}          & \textbf{120\%}       & \textbf{173\%} \\
  \hline
  \texttt{blake2b\_256} builtin, 20 levels & 7{,}974{,}399 & 16{,}254 \\
  \textbf{Utilisation}          & \textbf{0.08\%}      & \textbf{0.10\%} \\
  \hline
\end{tabular}
\end{center}

A depth-20 Poseidon insertion exceeds the CPU budget by 20\% and the memory
budget by 73\%. It does not fit, and it does not nearly fit. Working backwards
from the memory budget, the deepest tree that fits is \textbf{11 levels}, an
anonymity set of 2{,}048, consuming essentially the entire transaction
budget, leaving nothing for the rest of the validator.

The same operation with the \texttt{blake2b\_256} builtin costs 0.08\% of the
CPU budget. \textbf{Poseidon on-chain is roughly 1{,}500 times more expensive
than blake2b on CPU and 1{,}750 times on memory.}

\begin{remark}
The Poseidon figures are estimates built from measured primitives, and they are
\emph{lower bounds}. They count only field multiplications and additions. A real
Aiken implementation also pays for traversing the round-constant lists, indexing
the MDS matrix, and constructing intermediate lists; all at the same
$\approx 2\times10^5$ steps per operation. The true cost is higher, not lower.

This is also precisely why \texttt{aiken-zk}'s \texttt{merkle\_tree\_checker}
ships with a mocked hash function and a ``do not use in production'' warning
(Section~\ref{sec:ch7-aikenzk}). It is not an oversight. There is nothing to
put there.
\end{remark}

%--------------------------------------------------------------------------
\section{Wall Two: The Hash Function Dilemma} \label{sec:ch8-wall2}

The obvious response to Wall One is: use blake2b then. It is a builtin, it is
1{,}500 times cheaper on-chain, problem solved.

It is not solved, because the hash appears on \emph{both} sides. The Merkle path
is recomputed on-chain during deposit \emph{and} inside the circuit during
withdrawal. And the two environments have exactly opposite preferences.

\begin{center}
\begin{tabular}{lrr}
  \hline
  \textbf{Hash (2-to-1)} & \textbf{On-chain (CPU)} & \textbf{In-circuit} \\
  \hline
  Poseidon      & 600{,}802{,}860 & $\approx$ 240 constraints \\
  \texttt{sha2\_256}     & 581{,}923       & $\approx$ 30{,}000 constraints \\
  \texttt{blake2b\_256}  & 398{,}720       & not in circomlib \\
  \hline
\end{tabular}
\end{center}

\begin{highlight}
\textbf{The hash functions Cardano computes cheaply are the ones circuits prove
expensively, and vice versa.} Poseidon is 125 times cheaper in-circuit and
1{,}500 times more expensive on-chain. There is no hash function that is cheap
on both sides of a Cardano ZK application today.
\end{highlight}

\subsection*{Following the SHA-256 Branch Honestly}

Suppose we accept the trade and build the tree with \texttt{sha2\_256}. What
happens? Something initially encouraging, thanks to a fact established in
Chapter~\ref{ch:offchain}:

\begin{outline}
Groth16 verification on Cardano costs 2{,}005{,}373{,}690 CPU steps
\emph{regardless of the circuit's size}. A 600{,}000-constraint circuit verifies
for exactly the same price as our 216-constraint password circuit.
\end{outline}

So on-chain, both sides now fit comfortably:

\begin{center}
\begin{tabular}{lrr}
  \hline
  \textbf{Transaction} & \textbf{CPU used} & \textbf{\% of budget} \\
  \hline
  Deposit: 20 $\times$ \texttt{sha2\_256} & 11{,}638{,}458 & 0.12\% \\
  Withdraw: 1 Groth16 verify + mint       & $\approx$ 2{,}050{,}000{,}000 & 20.5\% \\
  \hline
\end{tabular}
\end{center}

The on-chain problem disappears. The cost has not gone away, though; it has
moved to the prover, and there it becomes severe:

\begin{center}
\begin{tabular}{lrr}
  \hline
  \textbf{Circuit} & \textbf{Poseidon tree} & \textbf{SHA-256 tree} \\
  \hline
  Constraints (depth 20)   & $\approx$ 4{,}800 & $\approx$ 600{,}000 \\
  Powers of Tau needed     & $2^{13}$          & $2^{20}$ \\
  Proving key size         & $\approx$ 2\,MB   & multiple GB \\
  Browser proving          & under a second    & not viable \\
  \hline
\end{tabular}
\end{center}

A 600{,}000-constraint Groth16 circuit needs a $2^{20}$ ceremony, a proving key
measured in gigabytes, and tens of seconds to minutes of proving on a laptop.
Nobody is downloading that into a browser tab.

And so the failure closes into a circle. Recall the three prover architectures
from Section~\ref{sec:ch7-where}. Client-side proving is the only one that
delivers privacy against the operator. A gigabyte proving key forces you to
Architecture B, a server-side prover. But the server-side prover receives the
\texttt{nullifier} and \texttt{secret}, which is the entire note.

\begin{highlight}
A mixer whose proofs are generated by a server is not a mixer. The operator can
link every deposit to every withdrawal, which is precisely the linkage the
construction exists to break. Choosing an on-chain-cheap hash forces
server-side proving, and server-side proving destroys the privacy property that
justified the mixer.
\end{highlight}

That is Wall Two, and it is a genuine dilemma rather than an engineering task.
Note that it would evaporate the moment Cardano gained a Poseidon builtin: the
tree would be Poseidon on both sides, the circuit would be 4{,}800 constraints,
and the proof would be generated in a browser in under a second.

%--------------------------------------------------------------------------
\section{Wall Three: eUTxO Contention} \label{sec:ch8-wall3}

Suppose Wall One and Wall Two both vanish tomorrow: Cardano ships a Poseidon
builtin, the tree costs almost nothing on-chain, the circuit is small, proving
happens in the browser. The mixer still does not work, for a reason that has
nothing to do with cryptography.

\subsection*{The Merkle Root Is Single-Threaded State}

The pool's Merkle root lives in the datum of one UTxO. Every deposit must
consume that UTxO and produce its successor. Two depositors who build
transactions against the same pool UTxO produce two transactions spending the
same input, and the ledger will accept exactly one. The loser's transaction
is not delayed; it is invalid, and must be rebuilt against the new root and
resubmitted.

\begin{center}
\begin{tabular}{lll}
  \hline
  \textbf{Property} & \textbf{Ethereum} & \textbf{Cardano} \\
  \hline
  Tree state          & Contract storage      & One UTxO's datum \\
  Concurrent deposits & Ordered by the miner  & All but one invalid \\
  Deposits per block  & Many                  & \textbf{One} \\
  Block time          & $\approx$ 12 s        & $\approx$ 20 s \\
  Max deposit rate    & Hundreds/min          & \textbf{3/min} \\
  \hline
\end{tabular}
\end{center}

Three deposits per minute is the theoretical ceiling, achieved only if depositors
coordinate perfectly. Under real contention (independent users, no
coordination) the effective rate is far lower, and users experience it as
transactions that simply fail.

\subsection*{Why This Is Fatal Rather Than Annoying}

For most contended Cardano applications, this is a UX problem with known
mitigations. For a mixer it is fatal, because \textbf{the anonymity set
\emph{is} the throughput}. Privacy in a shielded pool comes from the number of
plausible alternative depositors. A pool that accepts three deposits a minute at
absolute best, and realistically far fewer, cannot accumulate a meaningful
anonymity set. The construction's security parameter and the platform's
bottleneck are the same quantity.

\subsection*{And the Standard Mitigations Do Not Apply}

\begin{itemize}
  \item \textbf{Batching.} A coordinator collects deposits and inserts them
        together. This works technically, and the coordinator learns exactly
        which commitments arrived from which addresses in which batch. The
        batcher becomes the deanonymiser.
  \item \textbf{Sharding into $N$ parallel pools.} Standard practice for
        Cardano concurrency, and here it divides the anonymity set by $N$. Ten
        pools give ten times the throughput and one tenth of the privacy.
        Anonymity does not shard.
  \item \textbf{Deferred insertion.} Deposits queue in their own UTxOs and are
        folded into the tree later. This works, but the queue UTxOs are public
        and sit there linking depositor addresses to commitments until they are
        absorbed, and the queue is itself a contended resource.
  \item \textbf{Hydra.} A Hydra head has no contention, and also a fixed,
        known participant set, which is the opposite of an anonymity set.
\end{itemize}

\begin{highlight}
Every standard fix for eUTxO contention works by introducing a coordinator or by
partitioning users. A mixer's security property is destroyed by coordinators and
by partitions. Wall Three is not a performance limitation; it is a structural
incompatibility between what makes eUTxO good (deterministic, local validation
with no shared mutable state) and what a shielded pool needs (high-throughput
shared mutable state).
\end{highlight}

%--------------------------------------------------------------------------
\section{The Smaller Obstacles} \label{sec:ch8-smaller}

For completeness, several further problems would need solving even if the three
walls came down.

\subsection*{Groth16's Per-Circuit Trusted Setup}

Groth16 requires a fresh Phase 2 ceremony for every circuit. Since the tree
depth is compiled into the circuit and each denomination is its own pool, every
(denomination, depth) pair needs its own multi-party ceremony, its own proving
key, and its own verification key. Five denominations means five ceremonies, and
each one is a trust assumption a user must evaluate.

This is the strongest practical argument for \textbf{PLONK} over Groth16. PLONK
uses a single \emph{universal, updatable} setup: one ceremony serves every
circuit up to a size bound, and it can be extended indefinitely by new
contributors. For an application family like this, that difference is decisive.

The obstacle is that no production-grade PLONK verifier exists in Aiken today.
Groth16 verification is three point decompressions and four pairings, which maps
neatly onto Cardano's builtins. A PLONK verifier needs polynomial commitment
openings, a transcript hash, and considerably more field arithmetic, and we
have just measured what field arithmetic costs in a Plutus script. The same
interpreter overhead that blocks on-chain Poseidon also stands between Cardano
and a practical PLONK verifier. We return to this in
Chapter~\ref{ch:future}.

\subsection*{Four Proofs Per Transaction}

At 2.005 billion CPU steps per verification against a 10 billion budget, a
transaction can carry at most four Groth16 verifications. Batched withdrawals
(an obvious efficiency measure, and a privacy measure, since batching
withdrawals blurs timing) cap out at four.

\subsection*{No Relayer Market}

Tornado's relayer network exists because relaying is profitable at Ethereum fee
levels. Cardano transaction fees are a fraction of an ADA, so the fee a relayer
could extract is small in absolute terms while the operational and legal risk is
not. Nothing prevents relayers from existing on Cardano; economics makes them
unlikely to.

\subsection*{Datum Growth and Sizes}

The pool datum holds 20 filled-node hashes plus bookkeeping, under a kilobyte,
comfortably within the 16{,}384-byte transaction limit. The withdrawal redeemer
carries a 192-byte proof. Sizes are not a binding constraint here, which is worth
saying explicitly: the limits that bite are execution units and contention, not
bytes.

%--------------------------------------------------------------------------
\section{Verdict} \label{sec:ch8-verdict}

\begin{center}
\begin{tabular}{lp{6.9cm}}
  \hline
  \textbf{Component} & \textbf{Status on Cardano today} \\
  \hline
  Withdrawal circuit
    & Works. Standard Circom over BLS12-381. \\
  \addlinespace
  On-chain proof verification
    & Works. 2.0B CPU, $\approx$ 0.15 ADA, constant in circuit size. \\
  \addlinespace
  Nullifier double-spend check
    & Works, and is \emph{better} than Ethereum's: enforced by the ledger
      via token uniqueness. \\
  \addlinespace
  Fixed denominations
    & Works. Trivial in eUTxO. \\
  \addlinespace
  Recipient/fee binding
    & Works. Public inputs the validator checks. \\
  \hline
  \addlinespace
  On-chain Merkle insertion with Poseidon
    & \textbf{Fails.} 120\% of CPU and 173\% of memory budget at depth 20. \\
  \addlinespace
  Merkle insertion with a builtin hash
    & On-chain fine; forces a 600k-constraint circuit, hence server-side
      proving, hence no privacy. \\
  \addlinespace
  Concurrent deposits
    & \textbf{Fails structurally.} One deposit per block per pool; anonymity
      set equals throughput. \\
  \hline
\end{tabular}
\end{center}

\begin{highlight}
\textbf{Tornado Cash cannot be built on Cardano today.} Not because Cardano
cannot verify zero-knowledge proofs (it verifies them cheaply and elegantly),
but because it cannot maintain ZK-friendly cryptographic state on-chain, and
because a shielded pool needs high-throughput shared mutable state that the
eUTxO model is deliberately designed not to provide.
\end{highlight}

The distinction is worth stating precisely, because it generalises well beyond
mixers:

\begin{outline}
Cardano today supports ZK applications where the \textbf{prover's statement is
about data the verifier already has}, and struggles with applications where the
statement is about \textbf{an accumulator the chain must maintain}.
\end{outline}

Our password lock in Chapter~\ref{ch:password} is the first kind: the hash sits
in the datum, the validator reads it, done. Janus Wallet's authentication is the
first kind. A mixer, a private rollup, an anonymous voting register with
on-chain membership updates; all the second kind.

%--------------------------------------------------------------------------
\section{What Would Have to Change} \label{sec:ch8-changes}

\begin{center}
\begin{tabular}{p{4.0cm}p{8.0cm}}
  \hline
  \textbf{Change} & \textbf{Effect} \\
  \hline
  A Poseidon builtin
    & Removes Walls One and Two entirely. Deposits become as cheap as
      blake2b, circuits stay at 4{,}800 constraints, proving returns to the
      browser. \emph{This single primitive is the highest-leverage change
      available.} \\
  \addlinespace
  Native scalar field arithmetic builtins
    & Removes the $\approx 2\times10^5$-steps-per-operation floor. Makes
      on-chain Poseidon, PLONK verifiers, and custom curve work viable.
      Partially addressed by CIP-109's \texttt{modularExponentiation}. \\
  \addlinespace
  A production PLONK verifier in Aiken
    & One universal ceremony instead of one per pool. Blocked today by the
      same field-arithmetic cost. \\
  \addlinespace
  Larger execution budgets
    & Helps at the margin only. The memory budget rose from 11.25M to 14M to
      16.5M, roughly 47\% in total. On-chain Poseidon at depth 20 needs
      about $17\times$ the current memory budget. Incremental parameter
      increases will not close a gap of that size. \\
  \addlinespace
  Contention-free shared state
    & The genuinely hard one. Requires either a coordinator (which breaks
      anonymity) or a different execution model. \\
  \hline
\end{tabular}
\end{center}

The last row is why Midnight exists. Cardano's answer to applications needing
private shared state is not to reshape the eUTxO ledger but to put them on a
partner chain designed for it, with selective disclosure built in, which is
also a deliberate answer to the regulatory question that opened this chapter.
Chapter~\ref{ch:future} takes that up.

%--------------------------------------------------------------------------
\section{Chapter Summary} \label{sec:ch8-summary}

\begin{itemize}
  \item Tornado Cash is three ingredients: a commitment Merkle tree, a nullifier
        set, and fixed denominations. The anonymity set is the pool's
        throughput.
  \item The withdrawal side ports cleanly to Cardano. The circuit is standard
        Circom, and verification costs 2.0B CPU steps regardless of circuit
        size.
  \item The nullifier set ports \emph{better} than on Ethereum: minting a token
        named after the nullifier hash makes double-spend prevention a ledger
        invariant rather than a script check.
  \item \textbf{Wall One:} no Poseidon builtin. Measured, the cost floor for any
        elementary operation in a Plutus script is $\approx 2\times10^5$ CPU
        steps and $\approx 600$ memory units, driven by interpreter overhead,
        not operand width, so no optimisation escapes it. A depth-20 Poseidon
        insertion needs 120\% of the CPU and 173\% of the memory budget.
  \item \textbf{Wall Two:} the hashes Cardano computes cheaply cost
        $\approx 30{,}000$ constraints in-circuit; the hash that costs
        $\approx 240$ constraints cannot be computed on-chain. Choosing the
        on-chain-cheap hash forces a 600k-constraint circuit, hence server-side
        proving, hence no privacy against the operator.
  \item \textbf{Wall Three:} the Merkle root is one UTxO. One deposit per block
        per pool, $\approx$ 3/minute at best. Every standard mitigation
        introduces a coordinator or partitions users, and both destroy anonymity.
  \item The general lesson: Cardano supports ZK where the statement concerns
        data the chain already holds, and struggles where the chain must
        maintain a ZK-friendly accumulator.
  \item A single Poseidon builtin would remove two of the three walls.
\end{itemize}
